\section{Protocol}\label{sec:protocol}

\subsection{Three main arms and a CTIS-only injection control}

A0 and A1 use the same per-class slice budget and are evaluated on the same sealed panel.
A2 approximates that budget through its random partition; C1 is an additional CTIS-only control,
not a fourth arm run across all five libraries.

\begin{description}
\item[A0, control.] For each class, sample a realised budget $m_c$ close to the planned
$0.60\,S_c$ slices from $r_c^{C}$ distinct recordings in C, the training partition, which contains no slice of a test recording. Nothing the model sees
shares a recording with anything it is scored on.

\item[A1, homologous exposure.] Same $m_c$ and same C-side recording count, but the entire
$Q^{\text{pool}}_c$ is injected and an equal number of C slices is removed. Much of this pool
shares recording identifiers with the sealed panel, but some pool recordings contribute no test
slices. Realised overlap is given in Supplementary Material S1; $\alpha_c$ is the full injected
pool fraction, not the fraction restricted to recordings represented in the test panel. In one
class, the suona technique cell C0296, four candidate pool slices repeat content that already sits
in the sealed panel and are removed by the duplicate-content rule before injection, so A1 trains
on four fewer slices than A0 and C1 (25{,}498 against 25{,}502) and the per-class budgets are equal
everywhere else.

\item[A2, slice-level reconstruction.] Draw the slice-level 80/10/10 partition of C $\cup$
$Q^{\text{pool}}$, i.e.\ a slice-level interpretation of the audited protocol, including same-recording
mixing. Training uses 80\% of that union; because the union is about $0.75\,S_c$, A2 trains on about
60\% of the original corpus. That shortfall is the price of matching the budget: the realised
training count lands within one percent of $m_c$ in every library, as Table~\ref{tab:a2} shows. A2 also
reports a self-test score on its own 10\% held-out share, which is an illustrative self-test score
under our reconstruction of a slice-level split over the pooled corpus; it is neither a
reproduction of the original authors' split nor the original benchmark score.
\item[C1, non-test injection (CTIS only).] Starting from A0, remove 4{,}177 slices and
replace them classwise with 1{,}665 slices from unused C content and 2{,}512 from R, all from
non-test recordings. Injection is disjoint from A0 samples, test samples, test recording IDs
and test content hashes. Realised counts equal A0 classwise, so the two arms
differ only in which slices are trained on; the total is 25{,}502, four more than A1's
25{,}498.
C1 uses seeds 0--2 and the same final-epoch recipe. Because it consumes R content, its
validation scores cannot be interpreted; it performs no validation-based selection.
This exploratory addition holds volume approximately fixed but not recording diversity.
\end{description}

A0 and A1 hold the number of training slices and the C-side recording count fixed, so a
difference between them cannot be attributed to more data or to more recording diversity within
C. What differs is whose content the model sees. A2 varies both mechanisms at once and therefore
serves as a plausibility check rather than as a third estimate of the same quantity.

\subsection{What the contrast estimates}\label{sec:estimand}

We state the estimand explicitly, because the arms differ in more than one respect.
The treatment condition replaces $|Q^{\text{pool}}_c|$ of the $m_c$ C slices with the whole
$Q^{\text{pool}}_c$ of the same class, keeping the budget at $m_c$; the control condition draws
$m_c$ slices from C.
Fixed across the two arms are the per-class slice count, the per-class count of C-side training
recordings, the epoch count and the resulting optimiser steps, as reported in Supplementary Material~S1, the sealed evaluation panel, the three training seeds and
the per-class label set. What the treatment changes is the number of pool recordings that reach
training, the identity of the content they carry, and the within-class diversity those recordings
add: on CTIS, 578 of 963 test recordings are represented in A1 training and none in A0.
We therefore call $A1-A0$ the \emph{pool-injection contrast} $\Delta_{\text{pool}}$ and read it as
the score change associated with injecting the held-out pool under an otherwise fixed budget.
It is not the effect of moving from a recording-level to a slice-level partition, because that
move also repartitions the pool; A2 is the arm that approximates such a move.
Table~\ref{tab:arms} summarises what each arm trains on.

\begin{table}[t]
\centering\footnotesize\setlength{\tabcolsep}{4pt}
\caption{What each arm trains on in the primary library (CTIS, 201 classes, 963 test
recordings). Recordings are distinct canonical identifiers contributing at least one training
slice. The per-class budget $m_c$ is identical in A0 and C1; A1 is four slices short in class C0296.}
\label{tab:arms}
\begin{tabular}{@{}llrrll@{}}\toprule
Arm & Training content & Slices & Recordings & Budget & Test-recording relation\\\midrule
A0 & C only & 25{,}502 & 3{,}017 & $m_c$ & disjoint, 0/963\\
A1 & C + all $Q^{\text{pool}}$, equal budget & 25{,}498 & 3{,}770 & $m_c$ & 578/963 shared\\
A2 & 80\% of $C\cup Q^{\text{pool}}$, slice level & 25{,}263 & 3{,}679 & $\approx m_c$ & same-recording mixing\\
C1 & C remainder + part of R & 25{,}502 & 3{,}684 & $m_c$ & non-test R, 0/963\\\bottomrule
\end{tabular}\end{table}

\subsection{Training and sealing}

The recipe is frozen: 224$\times$224 input, ImageNet initialisation, ten epochs, AdamW with
learning rate $3\times10^{-4}$ and weight decay 0.05, cosine schedule with one warm-up epoch,
label smoothing 0.1, batch size 64, and two backbones --- AlexNet \citep{alexnet2012} and
Swin-T \citep{swin2021} --- under three seeds.
No data augmentation is applied: the recipe specifies none, and time-axis flips are not
meaningful for spectrograms. The number of epochs is fixed at ten and the final-epoch weights are
the ones evaluated: no checkpoint, epoch or seed is chosen on any panel, so the reported scores
carry no selection step, and the validation partition is monitored and reported but never used to
pick a model.

The sealed panel is evaluated once per run, after training, and never during it. Per-seed
predictions are written to a staging directory and moved into a sealed directory only when all
seeds of a library--arm--backbone group have completed; the sealed record contains the
group's seed spread. This local record is not an independent timestamp and does not by itself
prove when analytical choices were made. No model, epoch or seed selection is performed; R contains no test
recording and is used for monitoring only, except that C1 consumes part of R as described above. This is what makes the reported differences interpretable: a
pipeline that peeked at the sealed panel to choose a checkpoint would produce a number that no
longer measures transfer.

\subsection{Metrics}

The primary metric is macro-F1 on the sealed panel. We compute it at two levels. At the slice
level, each slice is one prediction. At the recording level, the predicted probabilities of all
slices of a recording are averaged and the recording receives one label, which gives every
recording the same weight within a library and keeps a library with $\bar K = 48$ from dominating
one with $\bar K = 8.7$ through sheer slice count. Both levels are reported side by side; the
primary contrast is at the slice level, as specified in the internal analysis plan, and the recording level
serves as a robustness check, with any disagreement between the two flagged in the text. Macro-averaging is a reasonable summary when class support is uneven; we take the
motivating argument from the literature on precision--recall summaries under imbalance \citep{saito2015pr}, though that argument concerns two-class problems and we do not treat it as a general justification. Macro-F1 averages only over classes with non-zero support or
predictions in the evaluation panel. In the reported intervals the class set is recomputed within
each resampling replicate: a class that is absent from a replicate is dropped from that
replicate's mean rather than scored as zero. Supplementary Material S1 reports the variant in
which the class set is fixed once at the panel level and absent classes score zero; the point
estimates are identical and the intervals differ by less than half a point.

\subsection{Statistics}

Differences are paired at the slice level: for a given library and backbone, all arms are
scored on the same slices, which we verify by comparing the label and recording arrays of the
sealed records element by element. Intervals come from 2{,}000 bootstrap resamples of test
\emph{recordings}, with whole recordings resampled and the same resampling indices applied to
every arm and backbone. These intervals describe test-recording resampling uncertainty
conditional on the realised split and the three trained seeds; they do not include uncertainty
from new training seeds or a new train/test split. Concretely, let $f_{s,a,b}$ be the macro-F1 of seed $s$, arm $a$ and
backbone $b$ on the sealed panel. A reported arm score is the mean over seeds,
$\bar f_{a,b} = \frac{1}{3}\sum_s f_{s,a,b}$, and a reported effect is
$\hat\Delta_b = \frac{1}{3}\sum_s (f_{s,A1,b} - f_{s,A0,b})$, so the printed arm columns,
rounded to four decimals, need not subtract to the printed effect. Within each bootstrap
replicate the per-seed paired difference is recomputed from the stored per-slice predictions and
the replicate value is the mean over seeds; the interval is the 2.5--97.5 percentile range and
the tail proportion is a two-sided smaller-tail resampling measure obtained by inverting the
percentile interval, not from a null distribution. The original script clipped
zero-tail estimates to $1/B$; the reported analysis uses the add-one convention
$2(k+1)/(B+1)$ for the smaller-tail count $k$, capped at one, with $B=2{,}000$.
Its minimum raw value is $2/2{,}001$, not evidence of an exact tail probability. Because the class set is fixed by each dataset rather than sampled from a population of classes, these intervals describe resampling uncertainty over test recordings within the realised split and do not extend to unseen classes or to other libraries.
Seed dispersion is reported as the range over
seeds and is not tested: three seeds are too few to estimate between-seed variance reliably,
so the primary intervals treat them as fixed and seed dispersion is descriptive. A supplementary
class-stratified recording bootstrap and joint seed--recording bootstrap assess sensitivity;
resampling three seeds does not replace additional training or split realisations.

The internal plan designates a primary family of two tests: CTIS under AlexNet and CTIS under Swin-T, both for
A1$-$A0. There is no independent public timestamp verifying when the plan was fixed. We correct within that family by Holm's method at $\alpha = 0.05$. The smallest effect
of interest, or SESOI, is the three-point macro-F1 threshold recorded in the internal plan.
It is a study-specific practical decision threshold, not a significance cutoff and not an
externally validated minimum useful gain. Precision does not justify practical importance.
The other four libraries are additional, non-confirmatory applications of the same operation,
not replications of an identical task. No public preregistration was made; Supplementary
Material S1 distinguishes the internal plan, execution changes and post hoc analyses.

\subsection{Negative controls}

If the mechanism is homologous content, libraries whose recordings contribute one slice have
nothing to inject: $Q^{\text{pool}}_c = \sum_i \lfloor k_i/2 \rfloor = 0$ in every class, so the
exposure arm coincides with the control and A1$-$A0 is not identifiable. The only arm that can
then differ is A2, which also changes the random partition; we therefore report A0 and A2 with three
seeds on three collections --- chest\_falsetto with 1{,}280 recordings of one slice each,
erhu\_playing\_tech with 1{,}253, and GZ\_IsoTech with 2{,}824 recordings, 2{,}899 slices and 46 recordings
with two to four segments --- and read the contrasts as a pipeline sanity check rather than a test of the main intervention.
The expansion from the initial single-seed controls to three seeds is exploratory and is
listed in Supplementary Material S1.

Guzheng\_Tech99 is excluded from the control runs: every recording contributes one slice, and
the release carries per-note technique annotations whose per-recording dominant label is the same
technique \emph{boxian} in all 99 recordings, so no recording-level class layer exists. GZ\_IsoTech
is kept, and its exposure is not exactly zero: 46 multi-segment recordings put nine slices into
the pool under one random draw, an exposure of 0.005, against an expectation of 0.0065 under 400 draws with a 95th
percentile 0.0098. Its realised exposure is therefore small but not zero; it is a low-exposure rather than a
zero-exposure control.

\subsection{Cross-instrument experiment}
\label{sec:crossinstrument}

One library supports a second question: whether the technique classification of CTIS transfers
to instruments never seen in training. We froze the class universe by a single rule fixed before
training: a technique--instrument cell entered if it held at least two recordings and its
instrument contributed at least three recordings, and six content-type labels were excluded as
material rather than technique (scales, open strings, fundamental tones, single tones, excerpts
and instrument-specific techniques). That leaves 43 techniques over 1{,}388 recordings from 49
candidate classes over 4{,}037 recordings. We then grouped instruments --- rather than recordings --- into three folds, so that each
recording is predicted exactly once by a model that has not seen its instrument. That property
holds for the instrument-grouped and recording-grouped partitions, not for the random
slice-level partition. Because that
partition changes both recording overlap and instrument novelty, a recording-grouped reference
partition provides an operational separation of the two successive contrasts. We build it by grouping slices of the same
recording, and of the same duplicate-audio group, into one fold, then balancing fold sizes across
classes; all 43 classes keep at least one test recording per fold and 110 instruments, drawn from 135
groups after name merging and each with at least three recordings, appear on at least two folds. The labels here are distinct from the main arms:
R-slice denotes random slice-level folds; R-rec denotes recording-grouped folds with instrument
overlap; L-inst denotes instrument-withheld folds; and L-inst+Adv adds an adversarial instrument head on the penultimate
representation through a gradient-reversal layer \citep{ganin2014,ganin2015} with weight ramped to 0.02 over two warm-up epochs,
and an auxiliary loss weight of 0.2. All four partitions use two seeds and three folds. Because held-out instruments never occur in training, the
instrument-identity probe that diagnoses the mechanism is fitted on an 80/20 split of the
\emph{training} recordings: a closed-set probe cannot identify unseen instrument classes.
The probe is standardised logistic regression with $C=1$ and at most 200 iterations, on frozen
penultimate features: AlexNet's 4{,}096-dimensional final hidden representation or Swin-T's
pooled pre-head representation. Each side is capped at 3{,}000 slices. Its AUC is the unweighted
mean of one-versus-rest binary AUCs for fitted instrument classes with positives and negatives
in the probe test subset; its separate top-1 score is instrument accuracy, not technique macro-F1.
The probe split uses seed 0 and has no test-recording overlap between fitting and evaluation.

\subsection{The internal plan and what changed}

The design was written down before the main runs, but it was not publicly registered and no
independent timestamp records when it was fixed; we therefore describe it as the internal plan and
treat the two CTIS tests as its primary family rather than as a registered hypothesis. Several
choices changed during execution; Supplementary Material S1 lists each one. The CTIS class
universe narrowed from 219 to 201 classes after the five-recording rule and placeholder removal.
Early injection targets of 50\% and 30\% of the training slices were abandoned in favour of
injecting the entire pool, which is the arm reported here. Per-class budgets are capped where C content is
insufficient; both arms train for the same ten epochs and the final-epoch weights are evaluated,
so a capped budget is not an early-stopping decision. The negative-control libraries were
expanded from one seed to three. C1 was added after the three main arms and is therefore
exploratory: it draws on unused C and R content, and its validation scores cannot be interpreted.
Recording-level point estimates were corrected after the main runs from a mean weighted by slice
count to one vote per recording. The retained intervals already used that per-recording
weighting, so only the point estimates changed; the slice-level estimates are unaffected, and
Supplementary Material S1 reports both. The class-stratified and joint seed--recording bootstraps, the self-test table and
the comparison against published scores were also added after the main runs. Two further analyses are
post hoc and are reported as such: the split of the CTIS sealed panel into test recordings that do
and do not contribute slices to the injected pool, and the fixed-class-set sensitivity table. Neither
is part of the internal plan, and neither is used to support the primary conclusion.

\subsection{Use of AI-assisted tools}

The authors used Coze Agent (ByteDance) and the GPT 6.1 Sol model, accessed through the DeepSider
API, between September and November 2026 for literature search, language editing, and help with
coding and debugging. No text was produced without author review; the authors verified every
reported number against the released artifacts and take full responsibility for the content.
